\documentclass[10pt,a4paper]{amsart}
\usepackage[margin=19mm]{geometry}
\usepackage{amsmath,amssymb,mathtools}
\usepackage[scr=rsfs,cal=euler]{mathalfa}
\usepackage{xfrac}
\usepackage[unicode,hidelinks,bookmarks=false]{hyperref}
\newcommand{\ie}{i.e.\ }
\newcommand{\cf}{cf.\ }
\newcommand{\resp}{respectively\ }
\newcommand{\Subsection}{\subsection}
\providecommand{\nobreakdash}{\nobreak\hbox{-}}
\setlength{\emergencystretch}{2em}
\multlinegap0pt
\usepackage{bm}
\usepackage{bbm}
\usepackage{dsfont}
\usepackage{xspace}
\usepackage{cancel}
\usepackage{mathtools}
\usepackage{amsthm}
\makeatletter
\@ifundefined{theorem}{\newtheorem{theorem}{Theorem}}{}
\@ifundefined{lemma}{\newtheorem{lemma}[theorem]{Lemma}}{}
\makeatother
\title[A note on the Burnside problem for homeomorphism groups of m]{A note on the Burnside problem for homeomorphism groups of manifolds}
\author{Donggyun Seo}
\date{Source: arXiv:2604.15627v2 (CC BY 4.0)}
\begin{document}
\maketitle

\noindent\emph{Source and licence.} A note on the Burnside problem for homeomorphism groups of manifolds. Version arXiv:2604.15627v2 (CC BY 4.0), distributed under the Creative Commons Attribution 4.0 International licence, as recorded in the arXiv licence metadata for this exact version and shown on the abstract page https://arxiv.org/abs/2604.15627v2. Full rights notice: licenses/arxiv-2604.15627-CC-BY-4.0.txt.\par\smallskip

\noindent\textit{Editorial scope.} Counted unit: the Lemma whose source label is \texttt{lem:tits} in the article (source0.tex lines 267--270, source label \texttt{lem:tits}), delivered together with the article's own complete original proof of it (source0.tex lines 272--292, one \texttt{proof} environment of 21 lines). No mathematics is written, repaired or re-derived: statement and proof are the article's own text line for line. Required context retained uncounted: none. Every definition, display and label of those lines is retained unchanged. Omitted from this package and not redistributed: 1--266, 271--271, 293--445. Nothing inside a retained excerpt is omitted or altered: every retained body is delivered exactly as the article's lines are written, so no omission receipt of any kind accompanies this package. Citations: the retained excerpts carry no citation command at all, so no bibliography entry of the article is transcribed and no reference is redistributed. Reference adaptation: none was needed and none was made; every reference inside the delivered unit resolves inside it, so no unresolved reference or citation is produced. Printed slips kept literal: no printed word, symbol, hypothesis or number of the counted statement or of its proof was altered. Layout and class: the article's own class is \texttt{amsart} with packages including amsfonts, mathrsfs, amsthm, amsxtra, amssymb, amsmath, amscd, hyperref, inputenc, t1enc, eucal, indentfirst, graphicx, subfig; the wrapper is the lane's pinned \texttt{amsart} editorial wrapper at 10pt on A4, which restates the theorem-like environments the retained text needs and adds the article's own macro definitions that the retained text needs (none), with extra wrapper packages (bm, bbm, dsfont, xspace, cancel, mathtools). The delivered numbering is the unit's own. Assets: no retained span contains a figure environment, a graphics-inclusion command, a PDF-page inclusion, a table or a TikZ picture; no image file is copied into this package, the package declares no asset dependency and no image file is redistributed with it. Licence: version arXiv:2604.15627v2 of this article is distributed under the Creative Commons Attribution 4.0 International licence, as recorded in the arXiv licence metadata for this exact version and shown on the abstract page https://arxiv.org/abs/2604.15627v2. A full rights notice is delivered at licenses/arxiv-2604.15627-CC-BY-4.0.txt. Characters: every retained excerpt is pure ASCII, so the delivered engine, XeLaTeX with its default Latin Modern text font, reports no missing character.
\begin{lemma} \label{lem:tits}
Let $G$ be a group satisfying the Tits alternative, i.e., every finitely generated subgroup of $G$ is either virtually solvable or contains a free subgroup of rank $2$. 
Then every finitely generated periodic subgroup of $G$ is finite.
\end{lemma}

\begin{proof}
Let $K \leq G$ be a finitely generated periodic subgroup.
By the Tits alternative, $K$ is either virtually solvable or contains a free subgroup $F \cong F_2$.
The latter is impossible, since $F_2$ contains elements of infinite order, contradicting $K$ being periodic.
Hence $K$ is virtually solvable, so there exists a solvable subgroup $S \leq K$ of finite index.

It remains to show that $S$ is finite.
Since $[K:S] < \infty$ and $K$ is finitely generated, the Schreier transversal lemma implies that $S$ is finitely generated.
As a subgroup of $K$, the group $S$ is also periodic.
We show by induction on the derived length $d$ of $S$ that every finitely generated periodic solvable group is finite.

If $d = 1$, then $S$ is abelian, and by the classification of finitely generated abelian groups, $S \cong \mathbb{Z}^r \oplus T$ for some $r \geq 0$ and finite group $T$.
Since $S$ is periodic, we have $r = 0$, so $S = T$ is finite.

If $d \geq 2$, the abelianization $S/S'$ is a finitely generated periodic abelian group, hence finite by the case $d=1$.
Since $[S:S'] < \infty$ and $S$ is finitely generated, the Schreier transversal lemma implies that $S'$ is finitely generated.
Since $S'$ is a subgroup of $S$, it is periodic, and it is solvable of derived length $d-1$.
By the induction hypothesis, $S'$ is finite, and therefore $|S| = [S:S'] \cdot |S'| < \infty$.

Hence $S$ is finite, and consequently $|K| = [K:S] \cdot |S| < \infty$.
\end{proof}

\end{document}
